% Author: Alfredo Sánchez Alberca (asalber@ceu.es)

\chapter{Introduction to R and RKWard}\label{cha:introduction}

\section{Introduction}
The big computing power achieved by computers has made them powerful tools essentials for all those disciplines, such as
Statistics, that require managing large volumes of data.
Nowadays, nobody addresses a serious statistical study without the help of data analysis software. 

\emph{R} is a powerful programming language that includes a lot of functions for representing and analyzing data. 
It was developed by Robert Gentleman and Ross Ihaka at the university of Auckland in New Zealand, although now is
maintained by a large scientific community all around the world.

\begin{center}
\includegraphics[scale=0.5]{chapters/introduction/img/Rlogo}
\end{center} 

The advantages of R with respect to other data analysis software like SPSS, SAS, Matlab, Minitqab or Excel, are
multiple:
\begin{itemize}
\item It's open source software, what implies that it's free. 
It can be downloaded from the web \url{http://www.r-project.org/}.
\item It's multi-platform .
There are versions for the main platforms: Windows, Macintosh, Linux, etc. 
\item It's supported by a huge scientific community that uses the software as an standard for data analysis. 
\item It has thousands of packages for performing any type of data analysis and graphics representations, from the most
common to the most innovative and cutting edge statistical procedures that are not included in other software. 
The packages are organized and documented in the CRAN (Comprehensive R Archive Network) repository, from where they
can be downloaded for free. 
In Spain there is a mirror of this repository in the web \url{http://cran.es.r-project.org/}.
\item It's programmable, what means that the user can create easily their own functions or packages for specific data
analysis.
\item There are a lot of books, manuals and tutorials for free, that allow the learning even for advances uses
in all kind of disciplines like Mathematics, Physics, Chemistry, Biology, Psychology, Medicine, etc. 
\end{itemize}

By default the working environment of R is the command line, what means that the computations are done with commands
manually typed in a text window.
However, there exists some graphical user interfaces (GUI) that ease its use, particularly for newbies.
The GUI that we are going to use along these practices is \emph{RKWard}, developed by Thomas
Friedrichsmeier, and the package rk.Teaching developed by Alfredo Sánchez Alberca in the CEU San Pablo
University and specifically designed for teaching Statistics.

The main goal of this chapter is to introduce the student to the use of RKWard and rk.Teaching, showing him the basic
operations to enter and manipulate data.

\section{Installation}
\subsection{Installation of R}
\begin{description}
\item[Linux] In the Debian distribution and derivatives (Ubuntu, Kubuntu, etc.) is as easy as typing the following
command in the command prompt
\begin{lstlisting}
> sudo apt-get install r-base-html r-cran-rcmdr r-cran-rodbc r-doc-html r-recommended
\end{lstlisting}
\item[Windows] Download from \url{http://cran.es.r-project.org/bin/windows/base/release.htm} the installation program,
execute it and follow the instructions. 
\end{description}

\subsection{Intallation of RKWard and rk.Teaching}
The GUI RKWard can be downloaded from the web \url{http://rkward.sourceforge.net/} where there are instructions for
the different platforms.

For Windows it is recommended to install the full package that contains R, RKWard and rk.Teaching available in  
\url{http://aprendeconalf.es/rkteaching/}

By default, the installation of R includes the basic packages for the most common operations or analysis.
However, for other analysis is required to install additional packages, like for instance the package
\variable{rkTeaching}, that includes menus and dialogs for most of the statistical procedures teached in these
practices. 

To install the package rk.Teching, first you have to download it from the web
\url{http://asalber.github.io/rkTeaching/}, then start R or RKWard and enter the following command in the command
console.
\begin{lstlisting}
> setwd("path_to_download")
> install.packages("rk.Teaching",repos=NULL,dep=True)
\end{lstlisting}

For packages hosted in the CRAN, it's possible to install them from RKWard selecting the
menu \menu{Settings > Manage R packages and plugins}.
In the dialog displayed, select the \mtab{Install/Update/Remove R packages} tab, then search or select the desired
package and click the button \button{OK}.
After installing a package, to use it in a work session it must be loaded. 
For that, in the same menu, select the \mtab{Load/Unload R packages} tab, then select the package, click the button
\button{Load} and finally the button \button{OK}.


\section{Solved exercises}
\begin{enumerate}[leftmargin=*]
\item Create a data set with the data in the sample below and save it with name \variable{colesterol.rda}
\begin{center}
\begin{tabular}{|l|c|r|r|r|}
\hline
\multicolumn{1}{|c|}{Name} & \multicolumn{1}{c|}{Gender} & \multicolumn{1}{c|}{Weight} & \multicolumn{1}{c|}{Height} &
\multicolumn{1}{c|}{Cholesterol}\\
\hline
José Luis Martínez Izquierdo  & H &  85 & 179 & 182\\
Rosa Díaz Díaz & F & 65 & 173 & 232\\
Javier García Sánchez  & M & 71 & 181 & 191\\
Carmen López Pinzón & F &  65 & 170 & 200\\
Marisa López Collado & F &  51 & 158 & 148\\
Antonio Ruiz Cruz & M & 66 & 174 & 249\\
\hline
\end{tabular}
\end{center}

\begin{indication}
To create a data set:
\begin{enumerate}
\item Select the menu \menu{File > New > Dataset}.
\item In the dialog displayed enter the name \variable{cholesterol} and click the button \button{OK}.
\item In the data editor window, define a variable in every column giving a name for the variable in the \field{Name}
row, a type (Numeric, Factor, String or Logical) in the \field{Type} row, and in the case of a factor, defining its
levels in the \field{Levels} row.
\item After define the variable, enter the data of every variable in the sample in the corresponding column.
\end{enumerate}
To save the data set:
\begin{enumerate}
\item Select the menu \menu{Workspace > Save Workspace}.
\item In the dialog displayed give the name \variable{cholesterol.rda} to the file, select the folder where to save it
and click the button \button{Save}.
\end{enumerate}
\end{indication}

\item Define a new variable \variable{Age} with the ages of the individuals in the sample, between the variables
\variable{Name} and \variable{Gender}.
\begin{center}
\begin{tabular}{|l|r|}
\hline
\multicolumn{1}{|c|}{Name} & \multicolumn{1}{c|}{Age} \\
\hline
José Luis Martínez Izquierdo & 18 \\
Rosa Díaz Díaz & 32 \\
Javier García Sánchez & 24 \\
Carmen López Pinzón & 35 \\
Marisa López Collado & 46 \\
Antonio Ruiz Cruz & 68 \\
\hline
\end{tabular}
\end{center}

\begin{indication} 
\begin{enumerate}
\item Click the left tab \button{Workspace}.
\item In the workspace window double-click the data set \variable{cholesterol} to edit it.
\item In the data editor window, right-click the column header of the variable \variable{Gender} and select
\menu{Insert new variable left}. 
\item In the new empty column enter the name and type of the variable \variable{Age}, and enter the age of every
individual.
\end{enumerate}
\end{indication}

\item Insert the following data of a new individual:
\begin{quote}
Name: Cristóbal Campos Ruiz.\\
Age: 44 years.\\
Gender: Male.\\
Weight: 70 Kg.\\
Height: 178 cm.\\
Cholesterol: 220 mg/dl.
\end{quote}

\begin{indication}
\begin{enumerate}
\item Insert the data of the new individual in the first empty row.
\end{enumerate}
\end{indication}

\item Create a new variable with the body mass index of every individual using the formula
\[
\text{bmi} = \frac{\text{Weight (in Kg)}}{\text{Height (in m)}^2}
\]

\begin{indication}
\begin{enumerate}
\item Select the menu \menu{Teaching > Data > Compute variable}.
\item In the dialog displayed enter the formula to compute the body mass index in the field \field{Variable
computation}.
\item In the field \field{Save as} click the button \button{Change}.
\item In the dialog displayed select as a parent object the data set \variable{cholesterol} and click
the button \button{OK}.
\item Enter the name \variable{bmi} for the new variable and click the button \button{Submit}.
\end{enumerate}
\end{indication}

\item Recode the body mass index variable into a new categorical variable \variable{Obesity} according to the following
rules:
\begin{center}
\begin{tabular}{lcl}
Less than de $18.5$ & \rightarrow & Low weight\\
From $18.5$ to $24.5$ & \rightarrow & Healthy\\
From $24.5$ to $30$ & \rightarrow & Overweight\\
Greater than $30$  & \rightarrow & Obese
\end{tabular}
\end{center}

\begin{indication}
\begin{enumerate}
\item Select the el menu \menu{Teaching > Data > Variable recoding}.
\item In the dialog displayed insert the variable \variable{bmi} in the field \field{Variable to recode}.
\item Enter the recoding rules below in the field \field{Recoding rules}:
\begin{quote}
\lstinline{lo:18.5 = 1}\\
\lstinline{18.5:24.5 = 2}\\
\lstinline{24.5:30 = 3}\\
\lstinline{30:hi = 4}
\end{quote}
\item In the field \field{Save as} click the button \button{Change}.
\item In the dialog displayed select as a parent object the data set \variable{cholesterol} and click the button
\button{OK}.
\item Enter the name \variable{Obesity} for the new variable and click the button \button{Submit}.
\item In the data edition window, enter the levels for the \variable{Obesity} factor, setting the label ``Low weight''
for the first category, ``Healthy'' for the second one, ``Overweight'' for the third and ``Obese'' for the fourth.
\end{enumerate}
\end{indication}


\item Filter the data set to get a new data set with the data of males. 
\begin{indication}
\begin{enumerate}
\item Select the menu \menu{Teaching > Data > Data filtering}.
\item In the dialog displayed insert the data set \variable{cholesterol} in the field \field{Data set}.
\item Insert the expression \lstinline{gender=="M"} in the field \field{Selection condition}.
\item Enter the name \variable{cholesterol.males} for the new data set and click the button \button{Submit}.
\end{enumerate}
\end{indication}

\end{enumerate}

\section{Proposed exercises}
\begin{enumerate}[leftmargin=*]
\item The data set \variable{neonates} of the package \variable{rk.Teaching}, contains information about a
sample of 320 newborns that meet the normal gestation time in a hospital during one year.
Do the following operations:
\begin{enumerate}
\item Load the data set.
\begin{indication}
\begin{enumerate}
\item Click the \field{Workspace} tab and double-click the \variable{rk.Teaching} package to unfold the data sets that
it contains. 
\item Right-click the data set \variable{nenonates} and select the menu \menu{Copy to .GlobalEnv} to copy the data
set to the working environment.
\end{enumerate}
\end{indication}

\item Compute the variable \variable{APGAR.average} as the mean of the variables \variable{APGAR1} and
\variable{APGAR5}.
\item Recode the variable \variable{weight} into the factor \variable{weight.category} with two categories
corresponding to weights less than and greater than $2.5$ Kg.
\item Recode the variable \variable{APGAR1} into the factor \variable{APGAR.state} with three categories: depressed
(APGAR$\leq 3$), moderately depressed ($3<$APGAR$\leq 6$) and normal (APGAR$>6$).
\item Filter the data set to get a new data set with the neonates of non-smoking mothers with an APGAR score at 1
minute less than or equal to 3.
How many neonates are there?
\end{enumerate}
\end{enumerate}
